\section{Introduction}

Search-and-rescue operations are among the most demanding environments for autonomous robot teams. Time pressure, partial observability, dynamic hazards, and the life-critical consequences of assignment errors create conditions under which fully autonomous systems cannot yet be trusted without human oversight. At the same time, the cognitive load of supervising a fleet of robots executing many concurrent rescue jobs quickly overwhelms a single human operator. Large language models offer an appealing middle ground: by reasoning over natural-language task descriptions and issuing tool calls to robot-control APIs, an LLM can handle routine assignment decisions autonomously, escalating to the human only when it is uncertain or when constraints are violated.

The Model Context Protocol (MCP) has recently emerged as a standardised JSON-RPC 2.0 interface for connecting LLMs to external tools and data sources. Applied to robotics, MCP enables a language model to query robot status, assign jobs, and dispatch navigation commands through a uniform, schema-validated API—an architecture we have already demonstrated in a MuJoCo-based multi-robot rescue system. However, a critical gap remains: when a human operator needs to correct an LLM decision, existing systems—including our own prior work and recent frameworks such as HMCF—treat the override as an out-of-band natural-language message that is appended to the conversation history. The LLM reads the correction, replans, and moves on. Crucially, the system retains no structured record of why the correction was made, what class of error it addressed, or how future decisions should be constrained to prevent recurrence.

This paper makes the observation that human overrides in a multi-robot rescue context fall into three semantically distinct classes: (i) constraint overrides, which correct violations of physical or operational limits (e.g., battery, payload, hazard zone); (ii) priority overrides, which adjust the urgency ordering of jobs based on information the LLM did not have; and (iii) strategy overrides, which introduce spatial or coordination heuristics that should govern planning over a longer horizon. We argue that these three classes carry different information, should be represented differently in the LLM's context, and should have different persistence lifetimes. Encoding all three as free text treats them identically and discards their semantic structure.

We propose Structured Authority Injection (SAI), a lightweight extension to our MCP server that exposes three typed override tools—one per correction class. When the human operator selects an override type and fills in a structured form in the supervision interface, the corresponding MCP tool is called, and the resulting structured record is appended to a persistent authority context that is prepended to every subsequent LLM prompt for the duration of the mission. This design has two key properties: the override is a first-class protocol event (logged, schema-validated, and auditable), and its semantic content is preserved in a form the LLM can reliably attend to across many planning cycles.

We evaluate SAI against an unstructured baseline—identical architecture but overrides expressed as free-text prompt appends—in a controlled experiment using our MuJoCo rescue simulation with GPT-4o. The experiment is designed to be executable within three weeks: 10 fixed scenarios across 3 override types and 2 injection modes, each repeated 5 times, for 300 total simulation runs. Our primary metric is the Recurrence Rate (RR): the fraction of post-override decisions that repeat the same class of error as the one that triggered the override. Secondary metrics are mission success rate, override count per mission, and correction latency.

The contributions of this paper are as follows:

\begin{enumerate}
\item We introduce Structured Authority Injection (SAI), the first mechanism to promote human overrides in an LLM-driven multi-robot system to typed, schema-validated MCP tool calls, with persistent propagation across planning cycles.
\item We define a three-class taxonomy of rescue-domain overrides (constraint, priority, strategy) and design MCP tool schemas for each, together with a persistence model that governs how each class modifies the LLM's planning context.
\item We provide controlled empirical evidence that structured override injection significantly reduces error recurrence and reduces the burden on the human operator, compared to unstructured free-text correction.
\item We discuss the implications for human-robot collaboration system design, in particular the role of protocol-level structuring as a lever for improving both safety and autonomy.

\end{enumerate}

The remainder of the paper is organised as follows. Section 2 reviews related work. Section 3 describes the system architecture. Section 4 presents the SAI design. Section 5 details the experimental methodology. Section 6 reports results. Section 7 discusses implications and concludes.